\chapter{Kazakh-FEVER 3K Annotation Protocols and Schema}
\label{app:fever_annotation}

This appendix provides the formal annotation guidelines, verification criteria, and quality assurance procedures utilized in constructing the \textbf{Kazakh-FEVER 3K} benchmark described in Chapter \ref{chap:factual_verification}.

\section{Gold Evidence Identification Criteria}
\label{sec:app_fever_evidence_criteria}

Annotators were instructed to identify the minimal, sufficient set of encyclopedic sentences from the 250,000-passage knowledge corpus required to verify or refute an asserted claim:
\begin{enumerate}
    \item \textbf{Minimality:} An evidence set must not include redundant sentences that do not directly bear on the truth value of the claim.
    \item \textbf{Sufficiency:} The selected evidence sentences must jointly entail or refute the claim without requiring external, unstated common-sense assumptions beyond standard linguistic competence.
    \item \textbf{Single-Hop vs. Multi-Hop Evidence:} Where a claim combines two distinct facts, annotators were permitted to select up to three connected evidence sentences from the same or hyperlinked encyclopedia passages.
\end{enumerate}

\section{Three-Way NLI Labeling Guidelines}
\label{sec:app_fever_labeling}

Each claim-evidence pair was assigned one of three mutual exclusive labels:
\begin{itemize}
    \item \textbf{\textsc{Supported}:} The claim is directly and unambiguously entailed by the cited evidence. Paraphrases, morphological variations, and syntactic reorderings that preserve exact truth conditions are labeled as Supported.
    \item \textbf{\textsc{Refuted}:} The claim is directly contradicted by the evidence. Contradictions include direct negation, contradictory temporal dates, conflicting geographical locations, or incompatible entity attributes.
    \item \textbf{\textsc{Not Enough Info (NEI)}:} The cited evidence is insufficient to confirm or deny the claim. This includes speculative statements, unverified claims, or instances where the evidence mentions the subject but omits the asserted predicate.
\end{itemize}

\section{Adversarial Hard NEI Perturbation Guidelines}
\label{sec:app_fever_hard_nei_guidelines}

To prevent models from relying on shallow lexical overlap, 50\% of the NEI partition was curated under our adversarial \textsc{Hard NEI} protocol:
\begin{enumerate}
    \item \textbf{High Lexical Overlap Neutrality:} The candidate claim must share over 80\% of its root morphemes with the candidate evidence sentence, yet assert an unmentioned attribute or relation that remains unverified.
    \item \textbf{Morphological Polarity Perturbations:} Claims generated by flipping verbal negation suffixes (\textit{-ма/-ме}) while introducing minor lexical substitutions, ensuring that models cannot rely solely on the presence of negative particles.
    \item \textbf{Evidential and Modal Boundary Cases:} Transforming definitive historical statements into hearsay or evidential assertions (\textit{-ыпты/-іпті}) where the supporting text only validates the event's occurrence but not the hearsay stance.
\end{enumerate}

\section{Quality Control and Inter-Annotator Agreement}
\label{sec:app_fever_agreement}

To guarantee benchmark integrity, the annotation followed a multi-stage validation pipeline:
\begin{itemize}
    \item \textbf{Native Speaker Qualifications:} All annotators were native Kazakh speakers holding advanced degrees in linguistics, computer science, or journalism.
    \item \textbf{Double-Blind Annotation:} A random 20\% validation sample ($N=600$) was independently labeled by two separate annotators without knowledge of prior labels.
    \item \textbf{Agreement Metric:} Inter-annotator agreement achieved a Cohen's $\kappa = 0.84$, indicating strong consensus.
    \item \textbf{Adjudication Protocol:} Discrepant labels were reviewed and adjudicated by an expert linguistic panel led by Dr. Irina Ualiyeva.
\end{itemize}
